\chapter{Sprint 1 : Fondations}

\section{Introduction}
Le Sprint 1 a produit le squelette porteur de \nova : le schema de base
de donnees, la couche de securite multi-locataire, le flux
d'authentification, la plomberie du webhook WhatsApp, et le cerveau IA
qui transforme les messages clients en rendez-vous reserves. L'objectif
de ce sprint, convenu lors de la planification, etait d'atteindre un
etat ou une seule clinique pouvait etre creee, son numero de telephone
enregistre aupres de Meta, et un client reel pouvait reserver un
rendez-vous reel de bout en bout sans aucune intervention humaine.

Le chapitre est organise selon l'ordre dans lequel le travail a
effectivement ete realise : d'abord le schema, puis la RLS, puis
l'authentification, puis le webhook, puis l'IA. Chaque section se
termine par l'approche de verification utilisee lors du sprint review.

\section{Backlog du sprint}
Le Sprint 1 a couvert les exigences fondamentales BF01--BF08, BF21--BF24,
BF31--BF32 du Chapitre II. Concretement :
\begin{itemize}
  \item Mise en place du projet Supabase et application du schema de base.
  \item Configuration de l'authentification email + lien magique.
  \item Mise en place des politiques Row-Level Security pour chaque table.
  \item Construction de la fonction Edge \texttt{whatsapp-webhook}.
  \item Construction de la fonction Edge \texttt{classify-message}.
  \item Construction de la fonction Edge \texttt{send-reminders} pour les
        rappels 24h et 1h.
\end{itemize}

\section{Schema de la base de donnees}

\subsection{Tables centrales}
Le schema central, defini dans \texttt{schema.sql}, comporte sept tables
plus leurs index et contraintes de cle etrangere.

\paragraph{businesses.} Chaque locataire est une ligne dans cette table.
Une entreprise possede un nom, un type (clinique, laboratoire, medecin,
salon, garage, entreprise), un plan d'abonnement, des identifiants
WhatsApp optionnels, et une colonne JSONB \texttt{birthday\_config}. La
colonne \texttt{owner\_user\_id} est la cle etrangere vers
\texttt{auth.users} qui alimente toutes les politiques RLS.

\paragraph{users.} Cette table est la table \emph{equipe}, a ne pas
confondre avec \texttt{auth.users}. Elle stocke les praticiens (medecins,
mecaniciens, coiffeurs) qui travaillent pour une entreprise. Un membre
peut optionnellement avoir un \texttt{auth\_uid} qui le lie a une
identite Supabase Auth.

\paragraph{clients.} Les clients. Une ligne client est creee
automatiquement la premiere fois qu'un numero de telephone envoie un
message a l'un des canaux WhatsApp de l'entreprise. Le schema a debute
simple --- nom, telephone, email --- et s'est etoffe au fil des sprints
pour inclure \texttt{age}, \texttt{birthday}, \texttt{profile} (JSONB),
\texttt{last\_visit\_at}, \texttt{no\_show\_count},
\texttt{medications}, \texttt{allergies}, \texttt{conditions}.

\paragraph{services.} Ce que vend l'entreprise. Un service possede un
nom, une duree par defaut en minutes, un prix optionnel, un medecin
prefere optionnel, une categorie optionnelle.

\paragraph{appointments.} La table transactionnelle centrale. Cles
etrangeres vers \texttt{clients}, \texttt{services} et \texttt{users}.
Le statut est contraint a une liste check (\texttt{pending},
\texttt{confirmed}, \texttt{cancelled}, \texttt{completed},
\texttt{no\_show}). La source est aussi contrainte
(\texttt{manual}, \texttt{whatsapp\_ai}, \texttt{whatsapp\_voice},
\texttt{web}, \texttt{api}).

\paragraph{messages.} Chaque message WhatsApp entrant et sortant est
journalise ici. Cette table est la piste d'audit.

\paragraph{notifications.} Le flux de l'icone cloche dans le tableau
de bord. Cree par divers agents automatises pour faire remonter ce
que le proprietaire devrait savoir.

\subsection{Discipline de migration}
Plutot que de modifier directement \texttt{schema.sql} a chaque ajout
de fonctionnalite, nous avons adopte la discipline d'ecrire un fichier
de migration distinct par feature pack : \texttt{enhancement\_pack.sql},
\texttt{lab\_pack.sql}, \texttt{packages\_pack.sql},
\texttt{pwa\_pack\_v2.sql}, \texttt{demand\_fill\_pack.sql},
\texttt{duration\_learning\_pack.sql}, \texttt{qr\_checkin\_pack.sql},
\texttt{drug\_check\_pack.sql}, \texttt{vehicle\_pack.sql},
\texttt{birthday\_pack.sql}, \texttt{queue\_pack.sql},
\texttt{voice\_avatar\_pack.sql}. Chaque pack est idempotent.

\section{Row-Level Security multi-locataire}
La multi-tenancy a ete la decision la plus consequente sur le plan
architectural. L'approche naive --- filtrer chaque requete SQL dans la
couche applicative par \texttt{business\_id} --- est fragile. Un seul
filtre oublie suffit a faire fuir les donnees entre locataires. La
Row-Level Security de Postgres resout ce probleme a la frontiere de la
base de donnees : chaque \texttt{select}, \texttt{insert},
\texttt{update} et \texttt{delete} passe par une politique qui
contraint automatiquement les lignes visibles a celles que
l'utilisateur authentifie courant est autorise a voir.

Une politique typique se presente ainsi :

\begin{lstlisting}[language=SQL,caption={Modele de politique RLS applique a la plupart des tables locataires}]
alter table public.appointments enable row level security;

drop policy if exists apt_owner on public.appointments;
create policy apt_owner on public.appointments
  for all using (
    business_id in (
      select id from public.businesses where owner_user_id = auth.uid()
    )
    or business_id in (
      select business_id from public.users where auth_uid = auth.uid()
    )
  );
\end{lstlisting}

La politique admet deux types d'utilisateurs : le proprietaire de
l'entreprise (dont \texttt{auth.uid()} correspond a
\texttt{businesses.owner\_user\_id}) et tout membre d'equipe lie via
\texttt{users.auth\_uid}. Les patients n'interagissent pas
directement avec cette table via PostgREST.

Nous avons applique le meme modele, avec des ajustements specifiques
par table, sur quinze tables.

\section{Authentification et onboarding}
L'authentification utilise le flux de lien magique de Supabase Auth.
Le proprietaire saisit son email sur la page d'accueil ; Supabase lui
envoie un lien a usage unique ; le clic ramene a l'application avec
un jeton d'acces dans le fragment d'URL. Le client JavaScript Supabase
le capte automatiquement, persiste la session dans \texttt{localStorage},
et le routeur du tableau de bord prend la suite.

Si l'utilisateur authentifie n'a pas encore d'entreprise, le routeur
le redirige vers \texttt{\#/create-business}, qui presente un
formulaire unique : nom de l'entreprise, type, telephone, adresse,
heures de travail.

Les gabarits d'email ont ete personnalises --- un email HTML de lien
magique avec le logo Nova et le gradient --- pour rendre la premiere
experience polie.

\section{Architecture du webhook WhatsApp}

\subsection{Verification Meta}
Avant qu'un webhook puisse recevoir des messages, Meta exige une
poignee de main de verification unique : Meta envoie un \texttt{GET}
avec \texttt{hub.verify\_token}, le webhook doit renvoyer
\texttt{hub.challenge} si le jeton correspond.

\begin{lstlisting}[language=JavaScript,caption={Poignee de main de verification webhook Meta}]
if (req.method === "GET") {
  const url = new URL(req.url);
  if (url.searchParams.get("hub.mode") === "subscribe"
      && url.searchParams.get("hub.verify_token") === VERIFY_TOKEN) {
    return new Response(url.searchParams.get("hub.challenge"));
  }
  return new Response("forbidden", { status: 403 });
}
\end{lstlisting}

\subsection{Routage des messages entrants}
La structure d'une charge utile Meta entrante est verbeuse. Chaque
callback contient un tableau \texttt{entry}, chaque entry contient un
tableau \texttt{changes}, chaque change contient un objet
\texttt{value} avec \texttt{messages} et \texttt{metadata}. Le
\texttt{display\_phone\_number} des metadonnees identifie quel de nos
numeros de telephone locataires a recu le message.

Une fois le locataire connu, nous extrayons le texte du message
entrant. Les messages texte simples nous donnent le corps directement.
Les notes vocales nous donnent un objet \texttt{voice} avec un media
id ; nous telechargeons le binaire depuis Meta, l'envoyons a OpenAI
Whisper pour transcription, et utilisons la transcription comme corps.
Les messages images --- pour les laboratoires recevant des photos de
prescriptions --- passent par le pipeline OCR Claude Vision decrit au
Chapitre IV.

\subsection{Idempotence et replays}
Meta ne garantit pas une livraison exactement-une-fois. Si notre
fonction repond lentement ou avec un statut non-200, Meta retentera le
meme callback. Pour eviter de doubler une reservation parce que Meta
a rejoue un message, le webhook deduplique contre la table
\texttt{messages} via un index unique sur \texttt{wa\_message\_id}.

\section{Le classifieur d'intention IA}

\subsection{Prompt systeme}
Le coeur de \texttt{classify-message} est le prompt systeme envoye a
Claude Haiku. Le prompt instruit le modele d'agir comme un
receptionniste virtuel pour un type d'entreprise specifique, dans la
langue detectee depuis le message entrant, et de repondre par un
objet JSON strict que l'executeur peut ensuite parser.

Le prompt complet fait environ 250 lignes et inclut le catalogue des
services disponibles pour ce locataire specifique (afin que Claude
n'invente pas des services qui n'existent pas), les heures de travail,
la liste FAQ, et une poignee d'indications de traduction
darija-vers-francais pour les mots que le modele lit parfois mal.

\subsection{Machine d'etats de prise de rendez-vous multi-tour}
Une vraie reservation prend plus d'un message. Le client dit \og je
veux un rendez-vous \fg{}, Nova demande le service, le client repond,
Nova demande la date, et ainsi de suite. Nous avons implemente cela
comme une machine d'etats adossee a une table \texttt{booking\_drafts}.

Lorsqu'un message entrant arrive, le cerveau verifie l'existence d'un
brouillon ouvert pour ce client. S'il en existe un, le message est
interprete dans le contexte de l'etat courant. Lorsque tous les champs
requis sont collectes, le cerveau insere le rendez-vous, supprime le
brouillon, et envoie la reponse de confirmation.

\subsection{Detection des conflits}
Avant d'inserer un rendez-vous, le cerveau interroge le calendrier
pour le creneau propose. Si le medecin est deja reserve a cette heure,
le cerveau redemande a Claude une alternative --- typiquement le
prochain creneau libre le meme jour ou le jour ouvrable suivant.

\subsection{Rappels}
La fonction Edge \texttt{send-reminders} s'execute toutes les quinze
minutes via \texttt{pg\_cron}. Elle recherche les rendez-vous prevus
dans les 24 prochaines heures qui n'ont pas encore eu leur rappel 24h
envoye, envoie un message WhatsApp interactif avec boutons Confirmer
et Annuler, et marque \texttt{reminder\_24h\_sent\_at}. Une seconde
passe gere les rappels 1h.

\section{Revue du sprint et tests}
A la fin du Sprint 1, nous avons execute un test fumee de bout en
bout : un projet Supabase frais, les onze packs SQL appliques, les
trois fonctions Edge deployees, un numero sandbox WhatsApp Business
Meta enregistre. Un client de test a envoye un message au numero
sandbox. En quatre secondes, Nova a repondu. Apres un echange de cinq
messages (salutation, service, date, heure, confirmation), le
rendez-vous est apparu dans la vue calendrier du tableau de bord. Le
rappel 24h s'est declenche le lendemain matin, le client a clique sur
Annuler, et le rendez-vous a ete marque annule.

La retrospective a fait remonter deux lecons. Premierement, les
limites de taux de l'API Meta sont plus strictes que ce que la
documentation suggere. Deuxiemement, Claude Haiku interprete parfois
\texttt{14h} comme \texttt{14:00} et parfois comme \texttt{02:00 PM} ;
nous avons resserre le prompt avec des exemples explicites pour
corriger.

\section{Conclusion}
Le Sprint 1 a livre le squelette porteur : une base de donnees
multi-locataire avec isolation RLS stricte, une authentification par
lien magique, un webhook WhatsApp qui gere la verification et
l'ingestion idempotente des messages, un cerveau IA qui classifie
l'intention et fait tourner une machine d'etats de prise de
rendez-vous multi-tour en trois langues, et un cron de rappels qui
ferme la boucle vingt-quatre heures plus tard. Le sprint suivant se
concentre sur les surfaces que le proprietaire utilise au quotidien :
le tableau de bord, le calendrier, le pipeline laboratoire, et la PWA
patient.
